\documentclass{article}
\usepackage[T1]{fontenc}
\usepackage{lmodern}
\usepackage{graphicx}
\usepackage{amsmath,amssymb}
\usepackage[a4paper,margin=2cm]{geometry}
\usepackage[absolute,overlay]{textpos}
\setlength{\TPHorizModule}{1mm}
\setlength{\TPVertModule}{1mm}
\usepackage[indonesian]{babel}
\usepackage{hyperref}
\setlength{\emergencystretch}{2em}
% unit: Halaman 27

\begin{document}

% === Halaman 27 (llm) ===

\section*{5. Replay Attack}

\begin{itemize}
    \item Tujuan: menggunakan kembali pesan sah yang pernah dikirim sebelumnya untuk menyebabkan suatu tindakan terjadi lagi.
    \item Yang menarik, pada \textit{replay attack} penyerang tidak perlu memecahkan enkripsi.
    \item Untuk mencegah \textit{replay attack}, biasanya digunakan mekanisme yang membuat setiap pesan mempunyai keunikan atau urutan, misalnya: \textit{Nonce, Timestamp}, atau \textit{Sequence number}
    \item Misalnya pesan diberi \textit{sequence number}:
    
    Message 101 \\
    Message 102 \\
    Message 103
    
    Jika penyerang mengirim kembali Message 101 setelah sistem sudah menerima Message 103, sistem dapat menolaknya karena nomor tersebut sudah kedaluwarsa.
\end{itemize}

\end{document}
